% Required packages in the preamble:
% \usepackage{longtable}
% \usepackage{array}
% \usepackage{booktabs}
% \usepackage{ragged2e}
% \newcolumntype{L}[1]{>{\RaggedRight\arraybackslash}p{#1}}

\appendix

\chapter{Complete Prompt Templates}
\label{app:prompts}

This appendix presents the complete prompt templates used in the experiments.
The placeholder \texttt{\{article\}} represents the processed article text
inserted at runtime. In the original S3 conditions, the retrieved examples
were selected dynamically using the article title as the BM25 query. The
methodological-basis column identifies research that informed each condition;
it does not imply that the prompt wording was copied from those studies.

\newcommand{\ArticleBlock}{Article: \texttt{\{article\}}\par Summary:}

\newcommand{\FramingDefinitions}{%
Left-leaning framing may include: emphasis on systemic inequality,
pro-government-intervention language, anti-corporate rhetoric, or progressive
social framing.\par
Centre/neutral framing: balanced reporting without clear partisan framing;
presents multiple perspectives or sticks to factual reporting without
editorial slant.\par
Right-leaning framing may include: emphasis on free-market solutions,
anti-regulation language, patriotic or traditional-values framing, or
anti-government-spending rhetoric.}

\newcommand{\NeutralDefinition}{%
Neutral, factual news reporting is characterised by presenting verifiable
facts without editorial judgment, covering multiple perspectives, and
avoiding charged or partisan language.}

\section{Original Prompt Conditions}
\label{app:original-prompts}

\footnotesize
\renewcommand{\arraystretch}{1.25}
\setlength{\tabcolsep}{4pt}

\begin{longtable}{@{}L{0.12\textwidth}L{0.61\textwidth}L{0.21\textwidth}@{}}
\caption{Original prompt conditions and their methodological basis}
\label{tab:original-prompts}\\
\toprule
Setting & Complete prompt & Methodological basis \\
\midrule
\endfirsthead
\multicolumn{3}{l}{\tablename\ \thetable\ continued}\\
\toprule
Setting & Complete prompt & Methodological basis \\
\midrule
\endhead
\midrule
\multicolumn{3}{r}{Continued on next page}\\
\endfoot
\bottomrule
\endlastfoot

S1 &
Summarize the following news article in 3 to 5 sentences.\par\medskip
\ArticleBlock
&
Baseline developed for this study.
\\
\midrule

S2 &
Summarize the following news article in 3 to 5 sentences. Your summary should
faithfully represent the original article's perspective and framing. Preserve
its political stance as written---do not amplify it, soften it, or shift it
toward neutrality.\par\medskip
\ArticleBlock
&
Perspective preservation \parencite{liu2024p3sum}.
\\
\midrule

S3a &
The following sentences represent neutral, factual reporting on a related
topic. Use them as a reference point to gauge the political spectrum. If the
article you are summarizing is more partisan than these sentences, your
summary should faithfully reflect that difference---do not move it closer to
neutral.\par\medskip
\texttt{\{three BM25-retrieved neutral sentences\}}\par\medskip
Now summarize the article below in 3 to 5 sentences.\par\medskip
\ArticleBlock
&
BM25 retrieval and external calibration data
\parencite{robertson2009probabilistic,kiesel2019semeval}.
\\
\midrule

S3b &
The following sentences illustrate how the same topic is reported across the
political spectrum. Use them to identify where the article you are summarizing
sits on this spectrum, then summarize it in 3 to 5 sentences while faithfully
preserving that position.\par\medskip
{[Left-leaning example]}\par
\texttt{- \{one BM25-retrieved left-leaning sentence\} [left-leaning]}\par\medskip
{[Neutral example]}\par
\texttt{- \{one BM25-retrieved neutral sentence\} [neutral]}\par\medskip
{[Right-leaning example]}\par
\texttt{- \{one BM25-retrieved right-leaning sentence\} [right-leaning]}\par\medskip
Now summarize the article below in 3 to 5 sentences.\par\medskip
\ArticleBlock
&
BM25 retrieval and spectrum calibration
\parencite{robertson2009probabilistic,kiesel2019semeval}.
\\
\midrule

S3c &
The following sentences illustrate how the same topic is reported across the
political spectrum. Use them to identify where the article sits on this
spectrum.\par\medskip
{[Left-leaning example]}\par
\texttt{- \{one BM25-retrieved left-leaning sentence\} [left-leaning]}\par\medskip
{[Neutral example]}\par
\texttt{- \{one BM25-retrieved neutral sentence\} [neutral]}\par\medskip
{[Right-leaning example]}\par
\texttt{- \{one BM25-retrieved right-leaning sentence\} [right-leaning]}\par\medskip
Summarize the article below in 3 to 5 sentences. Your summary should faithfully
represent the original article's perspective and framing. Preserve its
political stance as written---do not amplify it, soften it, or shift it toward
neutrality.\par\medskip
\ArticleBlock
&
Retrieval-based calibration and perspective preservation
\parencite{robertson2009probabilistic,kiesel2019semeval,liu2024p3sum}.
\\
\midrule

S4 &
\FramingDefinitions\par\medskip
Summarize the following news article in 3 to 5 sentences. Identify which of
the above framings (left / centre / right) best describes the article, then
write a summary that faithfully preserves that framing---do not amplify or
remove it. Output the summary only, with no preamble or framing analysis.\par\medskip
\ArticleBlock
&
Political framing and ideology labels
\parencite{entman1993framing,baly2020detect}. The definitions were
operationalised for this study.
\\
\midrule

S5 &
Read the article below and complete the two steps.\par\medskip
Step 1---Stance identification: In one sentence, describe the article's
political stance and key framing choices (e.g., which perspective it
emphasises, what language it uses).\par
Step 2---Faithful summary: Write a 3 to 5 sentence summary that faithfully
reflects that stance as identified in Step 1---do not amplify it or soften
it.\par\medskip
Format your response exactly as:\par
Stance: \texttt{<one sentence>}\par
Summary: \texttt{<3--5 sentences>}\par\medskip
Article: \texttt{\{article\}}
&
Step-by-step prompting and perspective preservation
\parencite{wei2022chain,liu2024p3sum}.
\\
\midrule

S6 &
Read the article and identify its political leaning. Then summarize it in 3
to 5 sentences while faithfully preserving that leaning.\par\medskip
\ArticleBlock
&
Perspective preservation \parencite{liu2024p3sum}.
\\

\end{longtable}

\section{Revised Prompt Conditions}
\label{app:revised-prompts}

The revised experiment retained S1 and introduced eight revised conditions.
These prompts were developed after the first-stage analysis. Their wording is
reported below exactly, apart from the replacement of runtime article text by
the \texttt{\{article\}} placeholder.

\begin{longtable}{@{}L{0.12\textwidth}L{0.61\textwidth}L{0.21\textwidth}@{}}
\caption{Revised prompt conditions and their methodological basis}
\label{tab:revised-prompts}\\
\toprule
Setting & Complete prompt & Methodological basis \\
\midrule
\endfirsthead
\multicolumn{3}{l}{\tablename\ \thetable\ continued}\\
\toprule
Setting & Complete prompt & Methodological basis \\
\midrule
\endhead
\midrule
\multicolumn{3}{r}{Continued on next page}\\
\endfoot
\bottomrule
\endlastfoot

S1 &
Summarize the following news article in 3 to 5 sentences.\par\medskip
\ArticleBlock
&
Unchanged baseline from the original experiment.
\\
\midrule

S2\textsubscript{r} &
Summarize the following news article in 3 to 5 sentences.\par
Your summary should faithfully represent the original article's perspective
and framing:\par
- If the article has a partisan stance (left or right leaning), preserve
it---do not soften or amplify it.\par
- If the article is neutral or centrist, write a neutral summary---do not
introduce any partisan framing.\par\medskip
\ArticleBlock
&
Symmetric revision developed from the first-stage findings; perspective
preservation \parencite{liu2024p3sum}.
\\
\midrule

S3a\textsubscript{r} &
\NeutralDefinition\par\medskip
Use this as a calibration baseline. If the article you are summarising departs
from neutral reporting---in a left or right direction---your summary should
faithfully reflect that departure. If the article is neutral, write a neutral
summary.\par\medskip
Now summarise the article below in 3 to 5 sentences.\par\medskip
\ArticleBlock
&
Definition-based calibration developed for this study; political framing
\parencite{entman1993framing}.
\\
\midrule

S3b\textsubscript{r} &
Political news reporting falls on a spectrum. Use the following as a
reference:\par\medskip
\FramingDefinitions\par\medskip
Identify where this article sits on the spectrum, then summarise it in 3 to 5
sentences while faithfully preserving that position. Do not shift it toward
or away from the centre.\par\medskip
\ArticleBlock
&
Definition-based spectrum calibration; framing and ideology labels
\parencite{entman1993framing,baly2020detect}.
\\
\midrule

S3c\textsubscript{r} &
Political news reporting falls on a spectrum. Use the following as a
reference:\par\medskip
\FramingDefinitions\par\medskip
Using this spectrum as a guide, summarise the article below in 3 to 5
sentences:\par
- If it is left or right leaning, preserve that framing faithfully---do not
soften or amplify it.\par
- If it is neutral or centrist, write a neutral summary---do not introduce
any partisan framing.\par\medskip
\ArticleBlock
&
Spectrum calibration and perspective preservation
\parencite{entman1993framing,baly2020detect,liu2024p3sum}.
\\
\midrule

S4\textsubscript{r} &
\FramingDefinitions\par\medskip
Read the article below. First, identify which framing (left / centre / right)
best describes it. Then write a 3 to 5 sentence summary that faithfully
preserves that framing---do not amplify or remove it.\par\medskip
Format your response as:\par
Framing: \texttt{<left / centre / right>}\par
Summary: \texttt{<3--5 sentences>}\par\medskip
Article: \texttt{\{article\}}
&
Explicit classification followed by generation; framing and perspective
preservation \parencite{entman1993framing,liu2024p3sum}.
\\
\midrule

S5\textsubscript{r} &
Read the article below and complete the two steps.\par\medskip
Step 1---Stance: Rate the article's political lean:\par
$-2$ = strongly left; $-1$ = leans left; $0$ = neutral/centre;
$+1$ = leans right; $+2$ = strongly right.\par
Then cite one specific phrase or framing choice from the article that
justifies your rating.\par\medskip
Step 2---Summary: Write a 3 to 5 sentence summary.\par
- If rating $\neq 0$: preserve the specific framing you cited---do not soften
or amplify it.\par
- If rating $=0$: write a factual, neutral summary---do not introduce any
partisan framing.\par\medskip
Format your response as:\par
Stance: \texttt{<rating> --- <one-sentence evidence>}\par
Summary: \texttt{<3--5 sentences>}\par\medskip
Article: \texttt{\{article\}}
&
Step-by-step prompting and evidence-grounded perspective preservation
\parencite{wei2022chain,liu2024p3sum}.
\\
\midrule

S5\textsubscript{adv} &
Read the article below and complete the three steps.\par\medskip
Step 1---Evidence: Quote 2 to 3 specific phrases or sentences from the article
that most clearly signal its political framing (word choice, emphasis, or
perspective). Use the article's exact words.\par\medskip
Step 2---Lean rating: Based ONLY on the evidence from Step 1, rate the overall
political lean:\par
$-2$ = strongly left; $-1$ = leans left; $0$ = neutral/centre;
$+1$ = leans right; $+2$ = strongly right.\par
In one sentence, explain how the quoted phrases support this rating.\par\medskip
Step 3---Summary: Write a 3 to 5 sentence summary.\par
- If lean $\neq 0$: the summary must reflect the framing signals identified
in Step 1---do not soften, amplify, or substitute with different partisan
framing.\par
- If lean $=0$: write a factual, balanced summary---do not introduce any
partisan framing.\par\medskip
Format your response as:\par
Evidence: \texttt{["<phrase 1>", "<phrase 2>", ...]}\par
Lean: \texttt{<rating> --- <one-sentence justification>}\par
Summary: \texttt{<3--5 sentences>}\par\medskip
Article: \texttt{\{article\}}
&
Evidence-first extension developed in this study, informed by step-by-step
prompting and perspective preservation
\parencite{wei2022chain,liu2024p3sum}.
\\
\midrule

S6\textsubscript{r} &
Read the article and identify its political leaning (left / centre / right).\par
Then summarise it in 3 to 5 sentences:\par
- If left or right leaning: preserve that framing faithfully---do not amplify
or soften it.\par
- If centre or neutral: write a factual, neutral summary---do not introduce
any partisan framing.\par\medskip
\ArticleBlock
&
Symmetric revision developed from the first-stage findings; perspective
preservation \parencite{liu2024p3sum}.
\\

\end{longtable}

\normalsize
